%%%%%%%%%%%%%%%%%%%%%%%%%%%%%%%%%
%
%       EXERCÍCIO
%
%%%%%%%%%%%%%%%%%%%%%%%%%%%%%%%%%

\ifdefstring{\atividade}{prova}{%
    \renewcommand{\valorquestao}{\ValorQ\ pontos}
}%

\begin{exercicioBanco}[\valorquestao]
O que acontece com a mediana, a média e o desvio-padrão de um conjunto de dados quando:
%
\begin{enumerate}[(a)]
    \item o valor 10 é somado a cada observação?
    %
    \item cada observação é multiplicada por 2?
    %
    \item subtrai-se a média de cada observação?
    %
    \item de cada observação, subtrai-se a média e divide-se pelo desvio-padrão?
\end{enumerate}
%
\end{exercicioBanco}

